%%=============================================================================
%% Voorwoord
%%=============================================================================

\chapter*{\IfLanguageName{dutch}{Woord vooraf}{Preface}}%
\label{ch:voorwoord}

%% TODO:
%% Het voorwoord is het enige deel van de bachelorproef waar je vanuit je
%% eigen standpunt (``ik-vorm'') mag schrijven. Je kan hier bv. motiveren
%% waarom jij het onderwerp wil bespreken.
%% Vergeet ook niet te bedanken wie je geholpen/gesteund/... heeft

Voor u ligt mijn bachelorproef, waarin ik onderzoek heb gedaan naar de mogelijkheden van een AI-gestuurd dashboard voor dependencybeheer. Het onderwerp sprak mij bijzonder aan omdat het een combinatie vormt van softwareontwikkeling, AI-intergatie en praktische probleemoplossing binnen een echte bedrijfscontext. Tijdens dit traject heb ik niet alleen mijn technische kennis verder verdiept, maar ook ervaren hoe belangrijk het is om een helder onderzoeksdoel te koppelen aan een bruikbare implementatie.

Deze bachelorproef was voor mij een leerzaam en uitdagend proces. Het uitwerken van de requitements, het ontwerpen van de architectuur, het bouwen van het proof-of-concept en het vergelijken van verschillende AI-aanpakken vroegen telkens om nieuwe inzichten en nauwkeurige keuzes. Vooral de vertaalslag van een conceptueel idee naar een werkende toepassing heb ik als bijzonder waardevol ervaren.

Graag wil ik enkele mensen bedanken die mij tijdens dit traject hebben ondersteund. In de eerste plaats dank ik mijn promotor, Thomas Aelbrecht, voor de begeleiding bij het schrijven en structureren van de paper, en voor de waardevolle feedback die mij geholpen heeft om mijn werk scherper en consistenter uit te werken. Daarnaast wil ik mijn co-promotor, Bert Pattyn, bedanken voor de inhoudelijke en praktische ondersteuning binnen Springbok zowel tijdens het papertraject als bij het uitwerken van de bachelorproef zelf. Ook gaat mijn dank uit naar Michiel Bogaert, die mij binnen Springbok heeft geholpen met het praktische gedeelte van het project en mij gerichte feedback gaf op mijn code en implementatie. Hun begeleiding en input hebben in belangrijke mate bijgedragen aan het resultaat van deze bachelorproef.

Tot slot wil ik ook mijn omgeving bedenken voor de steun en het geduld tijdens dit traject. Hun aanmoediging heeft mij geholpen om dit project met volharding af te ronden

\bigskip

\noindent
Jelle Vandriessche

% \lipsum[1-2]